\chapter{Parallelism on CPUs}\label{ch:cpupar}

WRF runs on CPU clusters with two levels of parallelism:
\begin{itemize}
  \item distributed memory, with MPI between processes;
  \item shared memory, with OpenMP threads inside a process.
\end{itemize}
This chapter describes both, because:
\begin{itemize}
  \item the CPU reference of the port runs on many MPI processes;
  \item the GPU version must give the same bits as one of them;
  \item the multi-GPU version of \cref{ch:mgpu} reuses the MPI structure.
\end{itemize}

\section{Patches and halos}

With MPI, the horizontal domain is split into a grid of \emph{patches}, \code{nproc_x} $\times$ \code{nproc_y}, one
per process. Every column belongs to exactly one patch. The vertical is never split. Each process allocates its
patch plus a \emph{halo}: a few rows and columns of its neighbours, wide enough for the widest stencil (five points
for fifth-order advection).

Before a computation that reads neighbours, the halo must hold the neighbours' current values. The exchanges are
declared in the Registry:
\begin{lstlisting}[language={}]
halo  HALO_EM_A  dyn_em  8:ru,rv,rw,ww,php,alt,al,p,muu,muv,mut,rho
\end{lstlisting}
This declares an exchange named \code{HALO_EM_A} of these fields with an 8-point stencil. The Registry turns it into
an include file that \code{solve_em} pulls in at the right place (\code{HALO_EM_A.inc}). The communication library is
RSL\_LITE (\code{external/RSL_LITE/}): it packs the halo rows of every field into a buffer, exchanges buffers with the
neighbours, and unpacks them.

\section{Tiles and threads}

Inside a patch, the work can be split again into \emph{tiles}, one per OpenMP thread (\code{numtiles}). The model
layer computes on tiles (\cref{ch:arch}). The tile loop is in \code{solve_em} and the physics drivers:
\begin{lstlisting}
!$OMP PARALLEL DO PRIVATE ( ij )
DO ij = 1 , grid%num_tiles
   CALL rk_tendency ( ..., grid%i_start(ij), grid%i_end(ij), &
                           grid%j_start(ij), grid%j_end(ij), ... )
END DO
\end{lstlisting}
Tiles share memory, so no halo exchange is needed between them. A routine must not write outside its tile.

\section{I/O quilting}

Writing history files from many processes can stall the computation. With quilting (\code{nio_tasks_per_group}),
extra MPI processes act as I/O servers: compute processes send their patches to them and continue, while the servers
assemble and write the file.

\section{Decomposition independence}

For a model to be trustworthy, its result should not depend on how many processes compute it. WRF is designed so that
it does not. Every point is computed by exactly one process with the same operations, and the halos supply the same
neighbour values that one process would have. There are two exceptions:
\begin{itemize}
  \item global reductions, such as a maximum or a sum over the domain, whose order can depend on the decomposition;
  \item loops whose bounds depend on the patch rather than the domain.
\end{itemize}
The port tests decomposition independence of the CPU reference before anything else (T-DEC: 1 vs 64 ranks, and 64 vs
128 vs 144 ranks, bitwise). A GPU run is one process. It can only equal a 64-process CPU run if the CPU result does
not depend on the number of processes.

\begin{portnote}
On one GPU there is one patch, with no MPI neighbours, and one tile. The halo exchanges become no-ops. The physical
boundary routines fill the domain-edge halo cells as on a single CPU process. For several GPUs (\cref{ch:mgpu}), the
same patch structure returns, with halos packed and exchanged directly from device memory.
\end{portnote}
